% Bagian: lambda-calculus
% Bab: lambda-definability
% Subbagian: lists

\documentclass[../../../include/open-logic-section]{subfiles}

\begin{document}

\olfileid[id]{lam}{ldf}{lst}
\olsection{Daftar}

Daftar $[M_1, M_2, \ldots, M_n]$ dapat didefinisikan sebagai berikut:
\[
  [M_1, M_2, \ldots, M_n] \ident \lambd[sz][s M_1 (s M_2 (\ldots (s M_n z)))]
\]
Secara informal, suku yang merepresentasikan suatu daftar merupakan fungsi
``akumulator'' daftar itu: fungsi yang menerima dua suku; suku pertama ialah
fungsi yang menerima nilai terakumulasi beserta suatu elemen baru dan
mengembalikan nilai terakumulasi yang baru; suku kedua ialah nilai
terakumulasi awal. Fungsi tersebut mengakumulasikan elemen satu per satu dan
pada akhirnya mengembalikan hasilnya.

\begin{digress}
  Jika Anda mengenal pemrograman fungsional, definisi ini akan tampak serupa
  dengan \emph{lipatan kanan}: suatu daftar didefinisikan
  sebagai fungsi lipatan kanan atas elemen-elemen yang dikandungnya.
\end{digress}

Kita dapat dengan mudah mendefinisikan beberapa fungsi berguna berdasarkan
pengodean ini:
\begin{align*}
  \fn{Sum} & \ident
  \lambd[l][l \, (\lambd[xa][\fn{Add} \, x \, a])\,  \num{0}]\\
  \fn{Len} & \ident
  \lambd[l][l \, (\lambd[xa][\fn{Add} \, a \, \num{1}]) \num{0}]
\end{align*}

$\fn{Sum}$ menghitung jumlah suatu daftar numeral Church. Fungsi ini bekerja
dengan melakukan akumulasi pada daftar tersebut: nilai awalnya ialah
$\num{0}$ dan, bagi setiap elemen, fungsi mengembalikan $\fn{Add} \, x\, a$
sebagai nilai baru (dengan $x$ sebagai elemen saat ini dan $a$ sebagai nilai
terakumulasi). Hasilnya ialah jumlah semua elemen.

\begin{prob}
  Apa yang dilakukan $\fn{Len}$? Jelaskan.
\end{prob}

\begin{prob}
  Definisikan fungsi yang membalikkan suatu daftar.
\end{prob}
\end{document}
